\section{Historical appendix: Scriptural Date and Location}
\label{historical-appendix}
% Generated. Do not edit.
% Deterministic projection of research/chronology.toml.
% schema: 5
% document: liturgy/roman-rite/1962/propers/temporal/f02-all-saints
% generated-by: tools/tpt proper-chronology annotations
\newcommand{\chronologyannotationclaim}[6]{#6}
\newcommand{\chronologyannotationreach}[3]{}
\newcommand{\chronologyannotationgroup}[3]{#3}
\newcommand{\chronologyannotation}[1]{%
  \ifcsname triptychchronologyannotation@#1\endcsname
    \csname triptychchronologyannotation@#1\endcsname
  \else
    \PackageError{triptych}{No chronology annotation for #1}{}%
  \fi}
% introit: Introit
\expandafter\def\csname triptychchronologyannotation@introit\endcsname{%
  \chronologyannotationgroup{traditional-attribution}{disputed}{\textbf{Traditional attribution} -- disputed: \chronologyannotationreach{Ps.32.1}{Ps.32}{true}\chronologyannotationclaim{israel.monarchy.david-traditional-era}{traditional-attribution}{catholic-traditional-v1}{disputed}{reigned from 1055 to 1015 B.C.}{David (reign in the usual chronology), B.C. 1055--1015}.}%
  \space \chronologyannotationgroup{historical-setting}{research-pending}{\textbf{Historical setting} -- Occasion date unresolved.}%
  \space \chronologyannotationgroup{composition}{preferred}{\textbf{Composition}: \chronologyannotationreach{Ps.32.1}{Ps}{true}\chronologyannotationclaim{critical.psalms.latest-composition-boundary}{composition}{catholic-critical-v1}{preferred}{before the Maccabean period, around 165 B.C.; no individual psalm can be dated securely}{Before c. 165 B.C.}}%
}
% lesson: Lesson
\expandafter\def\csname triptychchronologyannotation@lesson\endcsname{%
  \chronologyannotationgroup{composition}{disputed}{\textbf{Composition} -- disputed: \chronologyannotationreach{Apoc.7.10}{Apoc}{true}\chronologyannotationreach{Apoc.7.11}{Apoc}{true}\chronologyannotationreach{Apoc.7.12}{Apoc}{true}\chronologyannotationreach{Apoc.7.2}{Apoc}{true}\chronologyannotationreach{Apoc.7.3}{Apoc}{true}\chronologyannotationreach{Apoc.7.4}{Apoc}{true}\chronologyannotationreach{Apoc.7.5}{Apoc}{true}\chronologyannotationreach{Apoc.7.6}{Apoc}{true}\chronologyannotationreach{Apoc.7.7}{Apoc}{true}\chronologyannotationreach{Apoc.7.8}{Apoc}{true}\chronologyannotationreach{Apoc.7.9}{Apoc}{true}\chronologyannotationclaim{composition.apocalypse}{composition}{catholic-traditional-v1}{disputed}{the reign of Domitian ... A.D. 81-96}{A.D. 81--96}; \chronologyannotationreach{Apoc.7.10}{Apoc}{true}\chronologyannotationreach{Apoc.7.11}{Apoc}{true}\chronologyannotationreach{Apoc.7.12}{Apoc}{true}\chronologyannotationreach{Apoc.7.2}{Apoc}{true}\chronologyannotationreach{Apoc.7.3}{Apoc}{true}\chronologyannotationreach{Apoc.7.4}{Apoc}{true}\chronologyannotationreach{Apoc.7.5}{Apoc}{true}\chronologyannotationreach{Apoc.7.6}{Apoc}{true}\chronologyannotationreach{Apoc.7.7}{Apoc}{true}\chronologyannotationreach{Apoc.7.8}{Apoc}{true}\chronologyannotationreach{Apoc.7.9}{Apoc}{true}\chronologyannotationclaim{composition.apocalypse}{composition}{catholic-traditional-v1}{disputed}{from the year 90 to 100 (approximately)}{c. A.D. 90--100}.}%
}
% gradual: Gradual
\expandafter\def\csname triptychchronologyannotation@gradual\endcsname{%
  \chronologyannotationgroup{traditional-attribution}{disputed}{\textbf{Traditional attribution} -- disputed: \chronologyannotationreach{Ps.33.10}{Ps.33 Ps.39 Ps.70 Ps.85 Ps.94 Ps.97}{true}\chronologyannotationreach{Ps.33.11}{Ps.33 Ps.39 Ps.70 Ps.85 Ps.94 Ps.97}{true}\chronologyannotationclaim{israel.monarchy.david-traditional-era}{traditional-attribution}{catholic-traditional-v1}{disputed}{reigned from 1055 to 1015 B.C.}{David (reign in the usual chronology), B.C. 1055--1015}.}%
  \space \chronologyannotationgroup{superscription-setting}{preferred}{\textbf{Superscription setting}: \chronologyannotationreach{Ps.33.10}{Ps.33 Ps.55}{true}\chronologyannotationreach{Ps.33.11}{Ps.33 Ps.55}{true}\chronologyannotationclaim{israel.monarchy.david-at-geth}{superscription-setting}{catholic-traditional-v1}{preferred}{A. M. 2944, A. C. 1060}{David's flight to Achis, king of Geth, A.M. 2944}.}%
  \space \chronologyannotationgroup{composition}{preferred}{\textbf{Composition}: \chronologyannotationreach{Ps.33.10}{Ps}{true}\chronologyannotationreach{Ps.33.11}{Ps}{true}\chronologyannotationclaim{critical.psalms.latest-composition-boundary}{composition}{catholic-critical-v1}{preferred}{before the Maccabean period, around 165 B.C.; no individual psalm can be dated securely}{Before c. 165 B.C.}}%
}
% alleluia: Alleluia
\expandafter\def\csname triptychchronologyannotation@alleluia\endcsname{%
  \chronologyannotationgroup{composition}{disputed}{\textbf{Composition} -- disputed: \chronologyannotationreach{Matt.11.28}{Matt}{true}\chronologyannotationclaim{composition.gospel-of-matthew}{composition}{catholic-traditional-v1}{disputed}{about A.D. 38-45}{c. A.D. 38--45}; \chronologyannotationreach{Matt.11.28}{Matt}{true}\chronologyannotationclaim{composition.gospel-of-matthew}{composition}{catholic-traditional-v1}{disputed}{about the year 40-42}{c. A.D. 40--42}; \chronologyannotationreach{Matt.11.28}{Matt}{true}\chronologyannotationclaim{composition.gospel-of-matthew}{composition}{catholic-traditional-v1}{disputed}{the years 40-45}{A.D. 40--45}; \chronologyannotationreach{Matt.11.28}{Matt}{true}\chronologyannotationclaim{composition.gospel-of-matthew}{composition}{catholic-traditional-v1}{disputed}{about the year 60-68}{c. A.D. 60--68}; \chronologyannotationreach{Matt.11.28}{Matt}{true}\chronologyannotationclaim{composition.gospel-of-matthew}{composition}{catholic-traditional-v1}{disputed}{about the years 64-67}{c. A.D. 64--67}; \chronologyannotationreach{Matt.11.28}{Matt}{true}\chronologyannotationclaim{composition.gospel-of-matthew}{composition}{catholic-traditional-v1}{disputed}{about the year 50}{c. A.D. 50}.}%
}
% gospel: Gospel
\expandafter\def\csname triptychchronologyannotation@gospel\endcsname{%
  \chronologyannotationgroup{narrated-event}{research-pending}{\textbf{Event} -- No narrated-event date in the chronology corpus.}%
  \space \chronologyannotationgroup{composition}{disputed}{\textbf{Composition} -- disputed: \chronologyannotationreach{Matt.5.1}{Matt}{true}\chronologyannotationreach{Matt.5.10}{Matt}{true}\chronologyannotationreach{Matt.5.11}{Matt}{true}\chronologyannotationreach{Matt.5.12}{Matt}{true}\chronologyannotationreach{Matt.5.2}{Matt}{true}\chronologyannotationreach{Matt.5.3}{Matt}{true}\chronologyannotationreach{Matt.5.4}{Matt}{true}\chronologyannotationreach{Matt.5.5}{Matt}{true}\chronologyannotationreach{Matt.5.6}{Matt}{true}\chronologyannotationreach{Matt.5.7}{Matt}{true}\chronologyannotationreach{Matt.5.8}{Matt}{true}\chronologyannotationreach{Matt.5.9}{Matt}{true}\chronologyannotationclaim{composition.gospel-of-matthew}{composition}{catholic-traditional-v1}{disputed}{about A.D. 38-45}{c. A.D. 38--45}; \chronologyannotationreach{Matt.5.1}{Matt}{true}\chronologyannotationreach{Matt.5.10}{Matt}{true}\chronologyannotationreach{Matt.5.11}{Matt}{true}\chronologyannotationreach{Matt.5.12}{Matt}{true}\chronologyannotationreach{Matt.5.2}{Matt}{true}\chronologyannotationreach{Matt.5.3}{Matt}{true}\chronologyannotationreach{Matt.5.4}{Matt}{true}\chronologyannotationreach{Matt.5.5}{Matt}{true}\chronologyannotationreach{Matt.5.6}{Matt}{true}\chronologyannotationreach{Matt.5.7}{Matt}{true}\chronologyannotationreach{Matt.5.8}{Matt}{true}\chronologyannotationreach{Matt.5.9}{Matt}{true}\chronologyannotationclaim{composition.gospel-of-matthew}{composition}{catholic-traditional-v1}{disputed}{about the year 40-42}{c. A.D. 40--42}; \chronologyannotationreach{Matt.5.1}{Matt}{true}\chronologyannotationreach{Matt.5.10}{Matt}{true}\chronologyannotationreach{Matt.5.11}{Matt}{true}\chronologyannotationreach{Matt.5.12}{Matt}{true}\chronologyannotationreach{Matt.5.2}{Matt}{true}\chronologyannotationreach{Matt.5.3}{Matt}{true}\chronologyannotationreach{Matt.5.4}{Matt}{true}\chronologyannotationreach{Matt.5.5}{Matt}{true}\chronologyannotationreach{Matt.5.6}{Matt}{true}\chronologyannotationreach{Matt.5.7}{Matt}{true}\chronologyannotationreach{Matt.5.8}{Matt}{true}\chronologyannotationreach{Matt.5.9}{Matt}{true}\chronologyannotationclaim{composition.gospel-of-matthew}{composition}{catholic-traditional-v1}{disputed}{the years 40-45}{A.D. 40--45}; \chronologyannotationreach{Matt.5.1}{Matt}{true}\chronologyannotationreach{Matt.5.10}{Matt}{true}\chronologyannotationreach{Matt.5.11}{Matt}{true}\chronologyannotationreach{Matt.5.12}{Matt}{true}\chronologyannotationreach{Matt.5.2}{Matt}{true}\chronologyannotationreach{Matt.5.3}{Matt}{true}\chronologyannotationreach{Matt.5.4}{Matt}{true}\chronologyannotationreach{Matt.5.5}{Matt}{true}\chronologyannotationreach{Matt.5.6}{Matt}{true}\chronologyannotationreach{Matt.5.7}{Matt}{true}\chronologyannotationreach{Matt.5.8}{Matt}{true}\chronologyannotationreach{Matt.5.9}{Matt}{true}\chronologyannotationclaim{composition.gospel-of-matthew}{composition}{catholic-traditional-v1}{disputed}{about the year 60-68}{c. A.D. 60--68}; \chronologyannotationreach{Matt.5.1}{Matt}{true}\chronologyannotationreach{Matt.5.10}{Matt}{true}\chronologyannotationreach{Matt.5.11}{Matt}{true}\chronologyannotationreach{Matt.5.12}{Matt}{true}\chronologyannotationreach{Matt.5.2}{Matt}{true}\chronologyannotationreach{Matt.5.3}{Matt}{true}\chronologyannotationreach{Matt.5.4}{Matt}{true}\chronologyannotationreach{Matt.5.5}{Matt}{true}\chronologyannotationreach{Matt.5.6}{Matt}{true}\chronologyannotationreach{Matt.5.7}{Matt}{true}\chronologyannotationreach{Matt.5.8}{Matt}{true}\chronologyannotationreach{Matt.5.9}{Matt}{true}\chronologyannotationclaim{composition.gospel-of-matthew}{composition}{catholic-traditional-v1}{disputed}{about the years 64-67}{c. A.D. 64--67}; \chronologyannotationreach{Matt.5.1}{Matt}{true}\chronologyannotationreach{Matt.5.10}{Matt}{true}\chronologyannotationreach{Matt.5.11}{Matt}{true}\chronologyannotationreach{Matt.5.12}{Matt}{true}\chronologyannotationreach{Matt.5.2}{Matt}{true}\chronologyannotationreach{Matt.5.3}{Matt}{true}\chronologyannotationreach{Matt.5.4}{Matt}{true}\chronologyannotationreach{Matt.5.5}{Matt}{true}\chronologyannotationreach{Matt.5.6}{Matt}{true}\chronologyannotationreach{Matt.5.7}{Matt}{true}\chronologyannotationreach{Matt.5.8}{Matt}{true}\chronologyannotationreach{Matt.5.9}{Matt}{true}\chronologyannotationclaim{composition.gospel-of-matthew}{composition}{catholic-traditional-v1}{disputed}{about the year 50}{c. A.D. 50}.}%
}
% offertory: Offertory
\expandafter\def\csname triptychchronologyannotation@offertory\endcsname{%
  \chronologyannotationgroup{composition}{disputed}{\textbf{Composition} -- disputed: \chronologyannotationreach{Wis.3.1}{Wis}{true}\chronologyannotationreach{Wis.3.2}{Wis}{true}\chronologyannotationreach{Wis.3.3}{Wis}{true}\chronologyannotationclaim{composition.book-of-wisdom}{composition}{catholic-traditional-v1}{disputed}{the time of ... Ptolemy IV Philopator (221-204 B.C.)}{In or after the time of Ptolemy IV Philopator (221-204 B.C.)}; \chronologyannotationreach{Wis.3.1}{Wis}{true}\chronologyannotationreach{Wis.3.2}{Wis}{true}\chronologyannotationreach{Wis.3.3}{Wis}{true}\chronologyannotationclaim{composition.book-of-wisdom}{composition}{catholic-traditional-v1}{disputed}{the time of ... Ptolemy VII Physicon (145-117 B.C.)}{In or after the time of Ptolemy VII Physicon (145-117 B.C.)}.}%
}
% communion: Communion
\expandafter\def\csname triptychchronologyannotation@communion\endcsname{%
  \chronologyannotationgroup{composition}{disputed}{\textbf{Composition} -- disputed: \chronologyannotationreach{Matt.5.10}{Matt}{true}\chronologyannotationreach{Matt.5.8}{Matt}{true}\chronologyannotationreach{Matt.5.9}{Matt}{true}\chronologyannotationclaim{composition.gospel-of-matthew}{composition}{catholic-traditional-v1}{disputed}{about A.D. 38-45}{c. A.D. 38--45}; \chronologyannotationreach{Matt.5.10}{Matt}{true}\chronologyannotationreach{Matt.5.8}{Matt}{true}\chronologyannotationreach{Matt.5.9}{Matt}{true}\chronologyannotationclaim{composition.gospel-of-matthew}{composition}{catholic-traditional-v1}{disputed}{about the year 40-42}{c. A.D. 40--42}; \chronologyannotationreach{Matt.5.10}{Matt}{true}\chronologyannotationreach{Matt.5.8}{Matt}{true}\chronologyannotationreach{Matt.5.9}{Matt}{true}\chronologyannotationclaim{composition.gospel-of-matthew}{composition}{catholic-traditional-v1}{disputed}{the years 40-45}{A.D. 40--45}; \chronologyannotationreach{Matt.5.10}{Matt}{true}\chronologyannotationreach{Matt.5.8}{Matt}{true}\chronologyannotationreach{Matt.5.9}{Matt}{true}\chronologyannotationclaim{composition.gospel-of-matthew}{composition}{catholic-traditional-v1}{disputed}{about the year 60-68}{c. A.D. 60--68}; \chronologyannotationreach{Matt.5.10}{Matt}{true}\chronologyannotationreach{Matt.5.8}{Matt}{true}\chronologyannotationreach{Matt.5.9}{Matt}{true}\chronologyannotationclaim{composition.gospel-of-matthew}{composition}{catholic-traditional-v1}{disputed}{about the years 64-67}{c. A.D. 64--67}; \chronologyannotationreach{Matt.5.10}{Matt}{true}\chronologyannotationreach{Matt.5.8}{Matt}{true}\chronologyannotationreach{Matt.5.9}{Matt}{true}\chronologyannotationclaim{composition.gospel-of-matthew}{composition}{catholic-traditional-v1}{disputed}{about the year 50}{c. A.D. 50}.}%
}
% presentation: orientation-v2
\newcommand{\chronologyorientationrow}[2]{%
  \ifcsname triptychchronologyorientation@#1\endcsname
    \csname triptychchronologyorientation@#1\endcsname{#2}%
  \else
    \PackageError{triptych}{No chronology orientation row for #1}{}%
  \fi}
\expandafter\def\csname triptychchronologyorientation@row-e88e943a6faaa20b\endcsname#1{%
Introit: Psalms 32:1 (Vulgate numbering; Hebrew 33); Gradual: Psalms 33:10--11 (Vulgate numbering; Hebrew 34) & Traditional attribution: David's reign. #1 & Disputed: \chronologyannotationreach{Ps.32.1}{Ps.32}{true}\chronologyannotationreach{Ps.33.10}{Ps.33 Ps.39 Ps.70 Ps.85 Ps.94 Ps.97}{true}\chronologyannotationreach{Ps.33.11}{Ps.33 Ps.39 Ps.70 Ps.85 Ps.94 Ps.97}{true}\chronologyannotationclaim{israel.monarchy.david-traditional-era}{traditional-attribution}{catholic-traditional-v1}{disputed}{reigned from 1055 to 1015 B.C.}{B.C. 1055--1015}\\[.25em]%
}
\expandafter\def\csname triptychchronologyorientation@row-bbfc04633a7fe6e8\endcsname#1{%
Introit: Psalms 32:1 (Vulgate numbering; Hebrew 33) & Historical setting. #1 & Occasion date unresolved\\[.25em]%
}
\expandafter\def\csname triptychchronologyorientation@row-fb6656afd6f4b4f0\endcsname#1{%
Introit: Psalms 32:1 (Vulgate numbering; Hebrew 33); Gradual: Psalms 33:10--11 (Vulgate numbering; Hebrew 34) & Composition: Psalms. #1 & \chronologyannotationreach{Ps.32.1}{Ps}{true}\chronologyannotationreach{Ps.33.10}{Ps}{true}\chronologyannotationreach{Ps.33.11}{Ps}{true}\chronologyannotationclaim{critical.psalms.latest-composition-boundary}{composition}{catholic-critical-v1}{preferred}{before the Maccabean period, around 165 B.C.; no individual psalm can be dated securely}{Before c. 165 B.C.}\\[.25em]%
}
\expandafter\def\csname triptychchronologyorientation@row-d7a171509b4d3dc5\endcsname#1{%
Gradual: Psalms 33:10--11 (Vulgate numbering; Hebrew 34) & Superscription setting: David's flight to Achis, king of Geth. #1 & \chronologyannotationreach{Ps.33.10}{Ps.33 Ps.55}{true}\chronologyannotationreach{Ps.33.11}{Ps.33 Ps.55}{true}\chronologyannotationclaim{israel.monarchy.david-at-geth}{superscription-setting}{catholic-traditional-v1}{preferred}{A. M. 2944, A. C. 1060}{A.M. 2944 (Haydock edition reporting Ussher)}\\[.25em]%
}
\expandafter\def\csname triptychchronologyorientation@row-43dce930fad06090\endcsname#1{%
Offertory: Wisdom 3:1--3 & Composition: The Book of Wisdom. #1 & Disputed: \chronologyannotationreach{Wis.3.1}{Wis}{true}\chronologyannotationreach{Wis.3.2}{Wis}{true}\chronologyannotationreach{Wis.3.3}{Wis}{true}\chronologyannotationclaim{composition.book-of-wisdom}{composition}{catholic-traditional-v1}{disputed}{the time of ... Ptolemy IV Philopator (221-204 B.C.)}{In or after the time of Ptolemy IV Philopator (221-204 B.C.)}; \chronologyannotationreach{Wis.3.1}{Wis}{true}\chronologyannotationreach{Wis.3.2}{Wis}{true}\chronologyannotationreach{Wis.3.3}{Wis}{true}\chronologyannotationclaim{composition.book-of-wisdom}{composition}{catholic-traditional-v1}{disputed}{the time of ... Ptolemy VII Physicon (145-117 B.C.)}{In or after the time of Ptolemy VII Physicon (145-117 B.C.)}\\[.25em]%
}
\expandafter\def\csname triptychchronologyorientation@row-c1e1f88c6b954af2\endcsname#1{%
Gospel: Matthew 5:1--12 & Event. #1 & Date unresolved\\[.25em]%
}
\expandafter\def\csname triptychchronologyorientation@row-8b85083df79d6e95\endcsname#1{%
Alleluia: Matthew 11:28; Gospel: Matthew 5:1--12; Communion: Matthew 5:8--10 & Composition: The Gospel of St Matthew. #1 & Disputed: \chronologyannotationreach{Matt.11.28}{Matt}{true}\chronologyannotationreach{Matt.5.1}{Matt}{true}\chronologyannotationreach{Matt.5.10}{Matt}{true}\chronologyannotationreach{Matt.5.11}{Matt}{true}\chronologyannotationreach{Matt.5.12}{Matt}{true}\chronologyannotationreach{Matt.5.2}{Matt}{true}\chronologyannotationreach{Matt.5.3}{Matt}{true}\chronologyannotationreach{Matt.5.4}{Matt}{true}\chronologyannotationreach{Matt.5.5}{Matt}{true}\chronologyannotationreach{Matt.5.6}{Matt}{true}\chronologyannotationreach{Matt.5.7}{Matt}{true}\chronologyannotationreach{Matt.5.8}{Matt}{true}\chronologyannotationreach{Matt.5.9}{Matt}{true}\chronologyannotationclaim{composition.gospel-of-matthew}{composition}{catholic-traditional-v1}{disputed}{about A.D. 38-45}{c. A.D. 38--45}; \chronologyannotationreach{Matt.11.28}{Matt}{true}\chronologyannotationreach{Matt.5.1}{Matt}{true}\chronologyannotationreach{Matt.5.10}{Matt}{true}\chronologyannotationreach{Matt.5.11}{Matt}{true}\chronologyannotationreach{Matt.5.12}{Matt}{true}\chronologyannotationreach{Matt.5.2}{Matt}{true}\chronologyannotationreach{Matt.5.3}{Matt}{true}\chronologyannotationreach{Matt.5.4}{Matt}{true}\chronologyannotationreach{Matt.5.5}{Matt}{true}\chronologyannotationreach{Matt.5.6}{Matt}{true}\chronologyannotationreach{Matt.5.7}{Matt}{true}\chronologyannotationreach{Matt.5.8}{Matt}{true}\chronologyannotationreach{Matt.5.9}{Matt}{true}\chronologyannotationclaim{composition.gospel-of-matthew}{composition}{catholic-traditional-v1}{disputed}{about the year 40-42}{c. A.D. 40--42}; \chronologyannotationreach{Matt.11.28}{Matt}{true}\chronologyannotationreach{Matt.5.1}{Matt}{true}\chronologyannotationreach{Matt.5.10}{Matt}{true}\chronologyannotationreach{Matt.5.11}{Matt}{true}\chronologyannotationreach{Matt.5.12}{Matt}{true}\chronologyannotationreach{Matt.5.2}{Matt}{true}\chronologyannotationreach{Matt.5.3}{Matt}{true}\chronologyannotationreach{Matt.5.4}{Matt}{true}\chronologyannotationreach{Matt.5.5}{Matt}{true}\chronologyannotationreach{Matt.5.6}{Matt}{true}\chronologyannotationreach{Matt.5.7}{Matt}{true}\chronologyannotationreach{Matt.5.8}{Matt}{true}\chronologyannotationreach{Matt.5.9}{Matt}{true}\chronologyannotationclaim{composition.gospel-of-matthew}{composition}{catholic-traditional-v1}{disputed}{the years 40-45}{A.D. 40--45}; \chronologyannotationreach{Matt.11.28}{Matt}{true}\chronologyannotationreach{Matt.5.1}{Matt}{true}\chronologyannotationreach{Matt.5.10}{Matt}{true}\chronologyannotationreach{Matt.5.11}{Matt}{true}\chronologyannotationreach{Matt.5.12}{Matt}{true}\chronologyannotationreach{Matt.5.2}{Matt}{true}\chronologyannotationreach{Matt.5.3}{Matt}{true}\chronologyannotationreach{Matt.5.4}{Matt}{true}\chronologyannotationreach{Matt.5.5}{Matt}{true}\chronologyannotationreach{Matt.5.6}{Matt}{true}\chronologyannotationreach{Matt.5.7}{Matt}{true}\chronologyannotationreach{Matt.5.8}{Matt}{true}\chronologyannotationreach{Matt.5.9}{Matt}{true}\chronologyannotationclaim{composition.gospel-of-matthew}{composition}{catholic-traditional-v1}{disputed}{about the year 60-68}{c. A.D. 60--68}; \chronologyannotationreach{Matt.11.28}{Matt}{true}\chronologyannotationreach{Matt.5.1}{Matt}{true}\chronologyannotationreach{Matt.5.10}{Matt}{true}\chronologyannotationreach{Matt.5.11}{Matt}{true}\chronologyannotationreach{Matt.5.12}{Matt}{true}\chronologyannotationreach{Matt.5.2}{Matt}{true}\chronologyannotationreach{Matt.5.3}{Matt}{true}\chronologyannotationreach{Matt.5.4}{Matt}{true}\chronologyannotationreach{Matt.5.5}{Matt}{true}\chronologyannotationreach{Matt.5.6}{Matt}{true}\chronologyannotationreach{Matt.5.7}{Matt}{true}\chronologyannotationreach{Matt.5.8}{Matt}{true}\chronologyannotationreach{Matt.5.9}{Matt}{true}\chronologyannotationclaim{composition.gospel-of-matthew}{composition}{catholic-traditional-v1}{disputed}{about the years 64-67}{c. A.D. 64--67}; \chronologyannotationreach{Matt.11.28}{Matt}{true}\chronologyannotationreach{Matt.5.1}{Matt}{true}\chronologyannotationreach{Matt.5.10}{Matt}{true}\chronologyannotationreach{Matt.5.11}{Matt}{true}\chronologyannotationreach{Matt.5.12}{Matt}{true}\chronologyannotationreach{Matt.5.2}{Matt}{true}\chronologyannotationreach{Matt.5.3}{Matt}{true}\chronologyannotationreach{Matt.5.4}{Matt}{true}\chronologyannotationreach{Matt.5.5}{Matt}{true}\chronologyannotationreach{Matt.5.6}{Matt}{true}\chronologyannotationreach{Matt.5.7}{Matt}{true}\chronologyannotationreach{Matt.5.8}{Matt}{true}\chronologyannotationreach{Matt.5.9}{Matt}{true}\chronologyannotationclaim{composition.gospel-of-matthew}{composition}{catholic-traditional-v1}{disputed}{about the year 50}{c. A.D. 50}\\[.25em]%
}
\expandafter\def\csname triptychchronologyorientation@row-ca811ebc12efaaee\endcsname#1{%
Lesson: Revelation 7:2--12 & Composition: The Apocalypse of St John. #1 & Disputed: \chronologyannotationreach{Apoc.7.10}{Apoc}{true}\chronologyannotationreach{Apoc.7.11}{Apoc}{true}\chronologyannotationreach{Apoc.7.12}{Apoc}{true}\chronologyannotationreach{Apoc.7.2}{Apoc}{true}\chronologyannotationreach{Apoc.7.3}{Apoc}{true}\chronologyannotationreach{Apoc.7.4}{Apoc}{true}\chronologyannotationreach{Apoc.7.5}{Apoc}{true}\chronologyannotationreach{Apoc.7.6}{Apoc}{true}\chronologyannotationreach{Apoc.7.7}{Apoc}{true}\chronologyannotationreach{Apoc.7.8}{Apoc}{true}\chronologyannotationreach{Apoc.7.9}{Apoc}{true}\chronologyannotationclaim{composition.apocalypse}{composition}{catholic-traditional-v1}{disputed}{the reign of Domitian ... A.D. 81-96}{A.D. 81--96}; \chronologyannotationreach{Apoc.7.10}{Apoc}{true}\chronologyannotationreach{Apoc.7.11}{Apoc}{true}\chronologyannotationreach{Apoc.7.12}{Apoc}{true}\chronologyannotationreach{Apoc.7.2}{Apoc}{true}\chronologyannotationreach{Apoc.7.3}{Apoc}{true}\chronologyannotationreach{Apoc.7.4}{Apoc}{true}\chronologyannotationreach{Apoc.7.5}{Apoc}{true}\chronologyannotationreach{Apoc.7.6}{Apoc}{true}\chronologyannotationreach{Apoc.7.7}{Apoc}{true}\chronologyannotationreach{Apoc.7.8}{Apoc}{true}\chronologyannotationreach{Apoc.7.9}{Apoc}{true}\chronologyannotationclaim{composition.apocalypse}{composition}{catholic-traditional-v1}{disputed}{from the year 90 to 100 (approximately)}{c. A.D. 90--100}\\[.25em]%
}

The dates below distinguish traditional attribution, a remembered occasion,
the composition of a book, and the event narrated in the Gospel. The
Gradual appoints only the latter clause of its second verse; the Gospel
ends before the last clause of verse 12; the Wisdom chant is selected and
adapted. The whole-verse references in this chronological orientation
identify the passages queried without enlarging those liturgical selections.

\begin{chronologyorientationtable}
\chronologyorientationrow{row-e88e943a6faaa20b}{Davidic titles in the received Psalter; royal reference era, not a common occasion.}
\chronologyorientationrow{row-bbfc04633a7fe6e8}{No particular occasion or place is identified for the Introit verse.}
\chronologyorientationrow{row-fb6656afd6f4b4f0}{Collection-wide composition boundary; individual writing places uncertain.}
\chronologyorientationrow{row-d7a171509b4d3dc5}{David before accession, at Geth (Gath); inherited chronological marker, with its source qualification.}
\chronologyorientationrow{row-43dce930fad06090}{Alexandria in Gigot's account; either persecution horizon permits later publication.}
\chronologyorientationrow{row-c1e1f88c6b954af2}{Jesus teaches on an unnamed mountain after Galilean ministry (Matt.~4:23--5:1).}
\chronologyorientationrow{row-8b85083df79d6e95}{Palestinian original in the traditional account; Greek rendering distinguished. Speech sites are a separate question.}
\chronologyorientationrow{row-ca811ebc12efaaee}{John on Patmos; the visionary throne is not an earthly writing place.}
\end{chronologyorientationtable}

Fuller source qualifications: Psalms, p.~\pageref{hist-psalms};
Wisdom, p.~\pageref{hist-wisdom}; Matthew, p.~\pageref{hist-matthew};
Apocalypse, p.~\pageref{hist-apocalypse}.

\subsection{Psalms: attribution, flight, and composition}
\label{hist-psalms}
Both Psalms have Davidic titles in the Douay--Rheims/Challoner Psalter.
Psalm 32 supplies no particular occasion; Psalm 33 names the dismissal
associated with David's altered behavior. The reign shown above follows
Corbett's ``usual chronology'' in the historical \work{Catholic
Encyclopedia}. Corbett immediately records alternatives placing the
reign thirty to fifty years later. A royal era attached to an attributed
author does not date each attributed poem.\footnote{John Corbett,
``King David,'' \work{Catholic Encyclopedia}, IV (1908), opening paragraph;
Ps.~32 and 33 titles, Douay--Rheims/Challoner.}

The flight to Achis, king of Geth, is before accession and is separate
from that royal era. The received Psalm title names Achimelech; the
narrative in 1~Kings 21 names Achis. The Haydock edition's chronological
marker is inherited from the surrounding narrative, rather than freshly
printed beside the flight itself. Its attribution to Ussher is retained
in the table. The related Haydock note explicitly mentioning Usher is
at Ps.~70:1; that separate note identifies the chronology's source,
without supplying the occasion of Psalm 33.\footnote{Haydock,
\work{Douay--Rheims Bible with commentary}, Loreto/Feeney memorial
facsimile (2014), pp.~366, 368--369, artifact pp.~396, 398--399,
1~Kings 20--21; p.~740, artifact p.~776, Ps.~70:1. This is inspection
of a modern facsimile of historical commentary, not a newly established
independent date.}

The USCCB Psalms introduction supplies the displayed collection-wide
upper composition boundary and states the difficulty of securely dating
an individual psalm. It neither fixes one occasion for these two poems
nor establishes a shared place of composition. The traditional titles,
the flight narrative, and this composition judgment answer different
questions.\footnote{USCCB, \href{https://bible.usccb.org/bible/psalms/0}{
``Psalms: Introduction,''} dated source state inspected 25 October 2026;
the registered response is the 28 September 2026 artifact.}

\subsection{Wisdom: the author's city and the persecution in view}
\label{hist-wisdom}
Gigot distinguishes the book's Solomonic literary voice from its
actual unknown authorship. His argument for Alexandria appeals to
Greek language and thought, Egyptian circumstances, and the intended
Jewish audience amid surrounding idolatry. Alexandria is therefore
his proposed place of composition, rather than the named site of an
individual death in Wisdom 3.\footnote{Francis E. Gigot,
\href{https://www.newadvent.org/cathen/15666a.htm}{``Book of Wisdom,''}
\work{Catholic Encyclopedia}, XV (1912), authorship and ``Place and
date of composition.''}

The two reigns in the table identify alternative periods of persecution
which Gigot considers to be in the author's view. He expressly permits
publication after those rulers' deaths. The visible ``in or after''
qualification retains that possibility: neither reign's end is a
latest possible date of writing. The alternatives are the historical
source's proposals, not a claim that every later scholar accepts
either one.

\subsection{Matthew: the unnamed mountain and the written Gospel}
\label{hist-matthew}
Matthew 4:23--5:1 places Jesus' teaching after activity in Galilee.
The mountain is not named. The Alleluia's invitation belongs to the
wider Galilean narrative, without a named site for its delivery.
The table's unresolved event date is independent of the dates proposed
for writing the Gospel.

Jacquier reports several estimates in the historical \work{Catholic
Encyclopedia}'s account of Matthew's composition. His discussion
distinguishes proposed early dates for an original Hebrew or Aramaic
form from questions about the extant Greek Gospel. Durand similarly
proposes an Aramaic original and calls the time of its rendering into
Greek indefinite. The several estimates above retain their distinct
qualifications; averaging them would supply a date neither source
asserts.\footnote{E. Jacquier,
\href{https://www.newadvent.org/cathen/10057a.htm}{``Gospel of St.
Matthew,''} \work{Catholic Encyclopedia}, X (1911), ``Date and place
of composition''; Alfred Durand,
\href{https://www.newadvent.org/cathen/14530a.htm}{``The New Testament,''}
XIV (1912), origin and composition of the books.}

\subsection{Apocalypse: Patmos and the heavenly vision}
\label{hist-apocalypse}
Apocalypse 1:9 identifies John's setting as Patmos; 1:4 addresses
churches in Asia. The throne and worshippers in chapter 7 belong
to the vision, rather than to a second earthly address. Van den
Biesen reports the tradition placing John's exile near the end of
Domitian's reign. The table's broader reign span is a reference
period, not a claim that the book took that entire period to write.
Durand gives the other approximate range displayed above and
acknowledges a different earlier proposal among critics of his
period. The generated orientation preserves the returned alternatives
without inventing a new event date for the heavenly worship.
\footnote{Apoc.~1:4,9; Christian van den Biesen,
\href{https://www.newadvent.org/cathen/01594b.htm}{``Apocalypse,''}
\work{Catholic Encyclopedia}, I (1907), ``Time and place'';
Durand, ``The New Testament,'' Johannine writings paragraph.}
